%% Section 7: Long-Run Effects

\section{Longer-Run Effects}
\label{sec:longrun}

\subsection{Persistence in the NLPS: Round 11 (April 2024)}

The medium-run estimates in Tables~\ref{tab:firm_did}
and~\ref{tab:heterogeneity} show a 7.7-percentage-point enterprise survival
gap 14~months after the peak crunch, together with the 0.315 log-point sales
advantage reconciled in Section~\ref{sub:nlps_results}. This section examines
the compositional dimensions of this persistence using additional
household-level outcomes from Round~11.

Two mechanisms could sustain persistent divergence. First, enterprises that
closed or contracted during the crunch may have lost market relationships,
supplier networks, or key workers that are costly to rebuild, a \emph{scarring}
effect on firm-level human and social capital. Second, the crunch may have
accelerated a permanent reallocation of economic activity from cash-intensive
markets toward digitally-enabled channels, raising the long-run returns to
digital adoption and widening the productivity gap between high- and
low-fintech states.

Evidence from the household wage-employment panel (R5 to R11, column~1 of
Appendix Table~\ref{tab:additional_outcomes}) is informative about which
mechanism dominates. High-fintech states exhibit \emph{larger declines} in
formal wage employment at R11, with a coefficient of $-0.042$ ($p = 0.006$).
At first, this appears to contradict the enterprise survival and
employment results. The explanation lies in composition such that in high-fintech states,
surviving enterprises absorbed displaced wage workers into own-account
enterprise employment during the crunch. By R11, a higher share of the working
population in high-fintech states is engaged in enterprise self-employment
(consistent with the survival advantage shown in Table~\ref{tab:firm_did})
rather than formal wage work, a shift that shows up as a \emph{decline} in
the wage-employment indicator. In low-fintech states, enterprise exit was more
widespread, leaving fewer self-employment opportunities and pushing displaced
workers into inactivity rather than enterprise formation. This reallocation
pattern is more consistent with the \emph{reallocation} mechanism than pure
scarring. The crunch appears to have permanently redirected economic activity
toward enterprise self-employment in states with adequate digital
infrastructure, while low-fintech states simply contracted.

\subsection{Cross-Wave Evidence from the WBES (2014 vs. 2025)}
\label{sub:wbes_longrun}

The World Bank Enterprise Survey was fielded in Nigeria in 2014 (pre-shock)
and again in 2025 (approximately two years post-peak), providing a
cross-section comparison of formal-sector firm characteristics.
Table~\ref{tab:wbes_crosswave} reports zone-level means of the
finance-exclusion cash-dependence index in each wave, the change, and
the corresponding change in power-outage incidence.

\begin{table}[H]
\begin{threeparttable}
\caption{WBES Cross-Wave: Cash Dependence and Power Outages, 2014 vs.\ 2025}
\label{tab:wbes_crosswave}
\small
\begin{tabular}{l
  S[table-format=1.3]
  S[table-format=1.3]
  S[table-format=-1.3]
  S[table-format=1.3]
  S[table-format=1.3]
  S[table-format=-1.3]}
\toprule
 & \multicolumn{3}{c}{Cash-dependence index}
 & \multicolumn{3}{c}{Pct.\ with power outages} \\
\cmidrule(lr){2-4}\cmidrule(lr){5-7}
{Zone} & {2014} & {2025} & {$\Delta$}
       & {2014} & {2025} & {$\Delta$} \\
\midrule
North Central & 0.397 & 0.506 & 0.109 & 0.807 & 0.799 & -0.008 \\
North East    & 0.360 & 0.541 & 0.181 & 0.832 & 0.672 & -0.160 \\
North West    & 0.445 & 0.489 & 0.044 & 0.689 & 0.600 & -0.089 \\
South East    & 0.435 & 0.561 & 0.126 & 0.893 & 0.787 & -0.106 \\
South South   & 0.475 & 0.457 & -0.018 & 0.880 & 0.853 & -0.027 \\
South West    & 0.435 & 0.458 & 0.023 & 0.755 & 0.871 & 0.116 \\
\midrule
All zones     & 0.425 & 0.502 & 0.077 & 0.810 & 0.764 & -0.046 \\
\bottomrule
\end{tabular}
\begin{tablenotes}
\footnotesize
\item \textit{Notes:} Cash-dependence index is the mean across three
  indicators: (i)~no checking or savings account (\texttt{k6}),
  (ii)~no active loan or line of credit (\texttt{k8}), and (iii)~rates
  access to finance as a major or very severe obstacle (\texttt{k30}).
  The index ranges from~0 (fully finance-integrated) to~1 (fully
  excluded). Zone-level means are computed from firm-level data.
  The 2014 sample covers 2,676 firms in 19 states (mapped to geopolitical
  zones); the 2025 sample covers 1,043 firms in all six zones.
  The two waves are treated as independent cross-sections.
  South West includes Lagos, the most financially developed state in
  Nigeria; the \emph{increase} in power outages there (column 6)
  likely reflects the worsening of the Lagos electricity grid
  independent of the cash crunch.
\end{tablenotes}
\end{threeparttable}
\end{table}

The cross-wave comparison provides descriptive evidence consistent with a
reinforcing exclusion dynamic. Cash dependence increased most in North
Central ($+0.109$) and North East ($+0.181$), the zones with the lowest
pre-shock fintech adoption, suggesting firms without digital access during
the crunch had less incentive to invest in such infrastructure post-crunch,
perpetuating their exclusion. The South South ($-0.018$) and South West
($+0.023$) show minimal change, consistent with both zones' higher baseline
digital payment penetration insulating firms from the long-run footprint of
the crunch.

These patterns are descriptive and cannot be given a causal interpretation.
With only six geopolitical zones, no firm-level panel identifier across
waves, and potentially differential compositional change in the 2025 sample,
these comparisons motivate rather than confirm the persistence mechanism.
They are nonetheless consistent in direction with the \NLPS{} panel evidence
and with models of technology adoption under network externalities, where an
initial advantage in digital infrastructure compounds over time.
